% Generated draft: requires mathematical and source review.
\documentclass[11pt]{article}
\usepackage[margin=25mm]{geometry}
\usepackage{fontspec}
\setmainfont{Latin Modern Roman}
\usepackage{amsmath,amssymb,mathtools,unicode-math}
\setmathfont{Latin Modern Math}
\usepackage{enumitem,needspace,xurl,hyperref}
\setlength{\parindent}{0pt}
\setlength{\parskip}{5pt}
\setlength{\emergencystretch}{3em}
\newcommand{\N}{\mathbb{N}}
\newcommand{\Z}{\mathbb{Z}}
\newcommand{\Q}{\mathbb{Q}}
\newcommand{\R}{\mathbb{R}}
\newcommand{\C}{\mathbb{C}}
\newcommand{\F}{\mathbb{F}}
\newcommand{\A}{\mathbb{A}}
\newcommand{\PP}{\mathbb{P}}
\newcommand{\E}{\mathbb{E}}
\newcommand{\eps}{\varepsilon}
\newcommand{\dd}{\,\mathrm{d}}
\begin{document}
\section*{UnsolvedMath ID: 10900046}
\textbf{Source problem code:} catalogue entry
\subsection{Definitions and mathematical statement}
For an integer \(n\geq1\), begin with \(n\) tetrahedra and their \(4n\) triangular faces. A face-gluing outcome consists of a pairing of faces together with a simplicial identification of the three vertices in each paired pair. Let \(X_n\) be the quotient of the disjoint tetrahedra by those identifications. The source calls the resulting spaces pseudomanifolds and emphasizes that vertex links are not required to be spheres, being described instead as essentially random triangulated surfaces.

Maher asks to investigate the properties of these random quotients, allowing the nonspherical vertex-link regime rather than conditioning on ordinary \(3\)-manifolds. To make a probability statement precise, one must specify a distribution \(\mu_n\) on the face pairings and their gluing maps, and any validity or connectedness conditions. For instance, choosing labelled-face pairings and all vertex bijections uniformly is a possible model, but is not explicitly selected by the source; orientation-compatible choices, exclusion of self-gluings or conditioning give different models.

No particular property, statistic or asymptotic assertion is proposed in the excerpt. The precise conditions ensuring that all vertex links are genuine surfaces, rather than more singular quotients, and the intended definition of pseudomanifold also remain unstated. Once these choices are made, typical-property questions refer to the resulting laws of \(X_n\), potentially as \(n\to\infty\); neither a law nor a limiting theorem is invented here.
\subsection{Short English statement}
Study random quotients formed by pairing tetrahedron faces, allowing general vertex-link surfaces. Specify the probability law and admissible gluings before asking for typical or asymptotic properties.
\subsection{Sources}
\begin{itemize}
\item Maher, attributed item 6.2 in Delp, Hoffoss and Manning, \emph{Problems In Groups, Geometry, and Three-Manifolds} (2015), supplied PDF page 5. \url{https://arxiv.org/pdf/1512.04620#page=5}; original-status link \url{https://arxiv.org/abs/1512.04620}. Recorded access 2026-07-29; extraction source-tex checked against rendered PDF text, as reported in the record.
\item The supplied summary names Dunfield--Thurston, Maher, Dunfield--Hirsch and later random-triangulation/pseudomanifold work, without complete references or a verified complete characterization.
\item Local UnsolvedMath registry, \texttt{data/unsolvedmath/problems.json}, permanent ID \(10900046\), exact source code \texttt{catalogue entry}; supplied OPEN-TRIAGE status.
\end{itemize}
\end{document}
